
%%%%%%%%%%%%%%%%%%%%%%%%%%%%%%%%%%%%%%%%%%%%%%%%%%%%%%%%%%%%%
\subsubsection{模型建立}

% ==== 开篇引文 ====
针对问题一，本节将从...角度建立...模型。具体过程如下。

% ==== 步骤1：<标题> ====
\textbf{步骤1：...}

首先，...。设...为...，则有：
\begin{equation}
    ... = ...
\end{equation}
其中，$...$表示...。

将...代入...，可得：
\begin{equation}
    ... = ...
\end{equation}
进一步展开得：
\begin{equation}
    ... = ...
\end{equation}

% ==== 步骤2：<标题> ====
\textbf{步骤2：...}

根据...的性质，可以建立以下关系：
\begin{equation}
    ... = ...
\end{equation}
其中，$...$为...，$...$为...。

由式(1)和式(2)联立可得：
\begin{equation}
    ... = ...
\end{equation}

% ==== 步骤3：<标题> ====
\textbf{步骤3：...}

% ...（按需增加步骤，每步≥6个公式）...

% ==== 步骤4：求解算法 ====
\textbf{步骤4：...算法}

% 算法流程公式 + 参数表

\begin{table}[H]
  \centering
  \caption{算法参数设置}
  \begin{tabularx}{\textwidth}{CCCC}
    \toprule
    参数名 & 参数含义 & 取值 & 选取依据 \\
    \midrule
    ... & ... & ... & ... \\
    \bottomrule
  \end{tabularx}
  \label{tab:参数1}
\end{table}

% ==== 总结桥接 ====
综上所述，本节完成了...模型的建立，得到了...（目标函数/评价体系/预测模型），接下来将基于此进行求解。

% ============================================================
\subsubsection{模型求解}

% ==== 引文 ====
基于上述模型，本节利用...数据进行求解。求解环境为...，核心参数如表\ref{tab:参数1}所示。

% ==== 结果表格 ====
求解得到以下结果：

\begin{table}[H]
  \centering
  \caption{...结果}
  \begin{tabularx}{\textwidth}{CCCC}
    \toprule
    指标1 & 指标2 & 指标3 & 指标4 \\
    \midrule
    ... & ... & ... & ... \\
    \bottomrule
  \end{tabularx}
  \label{tab:结果1}
\end{table}

% ==== 结果图 ====
如图\ref{fig:xxx}所示，...（一句话点出图的结论）。

% 图引用示例：
% \begin{figure}[H]
%   \centering
%   \includegraphics[width=0.8\textwidth]{../求解/问题一/figures/xxx.png}
%   \caption{...}
%   \label{fig:xxx}
% \end{figure}

% ==== 结果分析 ====
由表\ref{tab:结果1}可知，...。从图\ref{fig:xxx}可以看出，...。
综合以上结果，...（一段分析总结）。
